\title{Kyllingesalat med karry}
\PrepTime{15}
\CookingTime{}
\CookingTempe{}
\TypeCooking{Pålægssalat}
\NbPerson{4}
\Image{0 100 800 900}{images/kyllingesalat.jpg}

\begin{ingredient}
    \begin{main}
        \item 400g Kylling\index[ingredient]{Kylling}
        \item 1,5dl Creme Fraiche 18\%\index[ingredient]{Creme Fraiche}
        \item 0,5dl Mayonnaise\index[ingredient]{Mayonnaise}
        \item 2 Forårsløg\index[ingredient]{Forårsløg}
        \item 1 tsk. Dijonsennep\index[ingredient]{Sennep}
        \item 1 tsk. Karry\index[ingredient]{Karry}
        \item 0,5 tsk. Løgpulver\index[ingredient]{Løgpulver}
        \item 0,25 tsk. Sukker\index[ingredient]{Sukker}
        \item Salt\index[ingredient]{Salt}
        \item Peber\index[ingredient]{Peber}
    \end{main}
\end{ingredient}
\begin{recipe}
    \step{Tilbered kyllingen på en pande.}
    \step{Rør creme fraiche, mayonnaise, dijon, karry og løgpulver sammen, og smag til med salt, peber og sukker.}
    \step{Vend kylling og forårsløg i dressingen og opbevar kyllingesalaten på køl til den skal serveres.}
\end{recipe}
\begin{notes}
    \textbf{Tip!} Man kan med fordel bruge lår kød for mere saftig kylling.
\end{notes}

\setcounter{step}{0}
\setcounter{substep}{0}

\clearpage